%% kapa: exam=1 round=24 year=2013 period=2 form=A subject=accounting subject_ko=회계학 seq=31 number=31
%% source: Q-Net 감정평가사 기출문제 (cst00309) artlSeq=5073802 file=extracted/1008/5073802/2013년 제24회 감정평가사 1차시험 기출문제/2교시(A형).hwp
%% conversion: kordoc 4.24.51 -> tools/gen1.py (auto); review=proofread-sonnet; flags=none
\begin{question}{31}
(주)감평은 두 개의 보조부문(수선부문과 동력부문)과 두 개의 제조부문(조립부문과 포장부문)으로 구성되어 있다. 수선부문에 집계된 부문원가는 노무시간을 기준으로 배부하며, 동력부문에 집계된 부문원가는 기계시간을 기준으로 배부한다. 보조부문원가를 제조부문에 배부하기 이전, 각 부문에 집계된 원가와 배부기준 내역은 다음과 같다. \par
{\small\centering
\begin{tabular}{|c|c|c|c|c|}
\hline
\multirow{2}{*}{\makecell{구분}} & \multicolumn{2}{|c|}{\makecell{보조부문}} & \multicolumn{2}{|c|}{\makecell{제조부문}} \\
\cline{2-5}
 & \makecell{수선부문} & \makecell{동력부문} & \makecell{조립부문} & \makecell{포장부문} \\
\hline
\makecell{노무시간} & \makecell{2,000시간} & \makecell{2,400시간} & \makecell{3,200시간} & \makecell{2,400시간} \\
\hline
\makecell{기계시간} & \makecell{5,000시간} & \makecell{5,000시간} & \makecell{10,000시간} & \makecell{10,000시간} \\
\hline
\makecell{부문원가} & \makecell{￦40,000} & \makecell{￦35,000} & \makecell{￦150,000} & \makecell{￦100,000} \\
\hline
\end{tabular}\par}

상호배부법을 사용하여 보조부문의 원가를 제조부문에 배부하면, 조립부문에 집계된 부문원가 합계액은? (단, 보조부문 용역의 자가소비분은 무시한다.) \par
\begin{choices}
  \item ￦135,000
  \item ￦185,000
  \item ￦190,000
  \item ￦195,000
  \item ￦200,000
\end{choices}
\end{question}
\answer{3}
